\section{Uniform height bounds for the complete collision source}
\label{sec:v44-complete-height}
The integrable power--logarithm class is larger than the class of
singularities which the physical roof density can actually have. We
prove a height bound before preparing any germ. The argument includes
arbitrarily small positive interior incidence and every section-decision
boundary. Its constants grow exponentially in the collision count;
it is not a local-limit estimate.

Let $\rho_{m,R}$ be the density of $L_{m,R}=\sum_{j<m}\tau_R\circ T_R^j$
under the full collision probability $\nu$. Put $R_-=9/20$, $R_+=47/100$,
$v_0=(3/50)/R_+=6/47$, and $c=\nu(Y_R^*)$. Densities and inequalities
between them are understood almost everywhere. We use full variation
mass, without the probability-distance factor $1/2$.

\subsection{Endpoint curvature without an interior incidence margin}
\begin{lemma}[The reduced action on an entire regular word]
\label{lem:v44-endpoint-curvature}
Fix a regular word of $m$ flights. The map from its initial state to
its two endpoint angles $(\alpha_0,\alpha_m)$ is injective and locally
invertible. On its open image the reduced action $F$ satisfies
\begin{equation}\label{eq:v44-endpoint-action}
 dF=R(-p_0\,d\alpha_0+p_m\,d\alpha_m),\qquad
 \nabla^2F\ge R\operatorname{diag}(c_0,c_m),
 \quad c_i=\sqrt{1-p_i^2}.
\end{equation}
Writing $B_m=\partial_p\alpha_m$ on the original branch, one has
\begin{equation}\label{eq:v44-endpoint-jacobian}
 |B_m|\ge v_0^m,\qquad
 d\alpha_0\,dp_0=|B_m|^{-1}d\alpha_0\,d\alpha_m.
\end{equation}
No assertion that the endpoint image is convex is needed.
\end{lemma}
\begin{proof}
Use oriented arclength at each contact. The second variation of a
segment is the sum of the two endpoint curvature terms
$c_j y_j^2/R+c_{j+1}y_{j+1}^2/R$ and the nonnegative transverse square
in the proof of Proposition~\ref{prop:general-edge}. Summing the
segments gives a positive definite contact Hessian. The internal
reflection equations are stationary equations, so the Hessian of the
reduced action is its internal Schur complement. Minimizing the full
quadratic form over the internal coordinates retains the two endpoint
curvature terms. Changing from arclength to angles gives the inequality
in \eqref{eq:v44-endpoint-action}. The first variation gives its equality.
All statements are local on a regular branch and remain valid as
interior incidence tends to zero through positive values.

For injectivity, fix both endpoint positions and minimize the polygonal
length over the product of the closed internal disks. At a regular
reflected contact its gradient is $-2c_j n_j$. For a different point
$q'_j$ of the same strictly convex disk,
$n_j\cdot(q'_j-q_j)<0$. The convex gradient inequality therefore makes
any different internal tuple have strictly larger length. Thus two
regular reflected paths of this word cannot share both endpoints.
For $m=1$ their straight segment determines the initial velocity directly.
This use of a convex minimization does not assert physical admissibility
of an arbitrary minimizing path: the word under consideration is already
physical, and other-disk exclusions are kept in its domain.

By \eqref{eq:positive-collision-matrix}, every entry of the positive
one-flight matrix $-DT_R$ is at least $v_0$. Every entry of the product
of $m$ such matrices is at least $v_0^m$. In particular $B_m\ne0$.
The change-of-variables formula now proves
\eqref{eq:v44-endpoint-jacobian}, with an absolute Jacobian for either parity.
\end{proof}

Use finitely many rational angular charts
$\alpha=\alpha_*+2\arctan x$, $|x|<1$, with a fixed semialgebraic disjoint
refinement at their overlaps. The chart derivative $a(x)$ lies in
$[1,2]$ and $|a'(x)|\le4$. Set $\delta_0=1/20$. On the part where
$|p_0|,|p_m|<\delta_0$, the endpoint-chart function $f(x,y)$ satisfies
\begin{equation}\label{eq:v44-chart-convexity}
 \nabla^2 f\ge\lambda_0 I_2,\qquad \lambda_0=R_-/2.
\end{equation}
Indeed the congruence of the first Hessian in
\eqref{eq:v44-endpoint-action} is at least
$R\sqrt{1-\delta_0^2}I_2$. The diagonal coordinate correction has norm
at most $4R\delta_0$. Since $\sqrt{1-\delta_0^2}-4\delta_0>1/2$,
\eqref{eq:v44-chart-convexity} follows. The collision-density weight
in these coordinates is bounded by $v_0^{-m}/\pi$.

\subsection{A coefficient-independent bound on level complexity}
\begin{lemma}[Length and turning of physical roof levels]
\label{lem:v44-level-complexity}
There are constants $C_1>1$ and $C_2>1$, depending only on the
one-flight geometry and the fixed charts, with the following properties.
For each fixed word and almost every roof value, its roof level in an
initial $(x,p)$ chart has total length at most $C_1C_2^m$.
On its near-normal endpoint-chart portions the total absolute curvature
of the regular roof level is at most $C_1C_2^m$.
These bounds are uniform in $R$ and require no distance from singular
or decision boundaries. They count every connected portion, not only
one selected inverse curve.
\end{lemma}
\begin{proof}
Keep the finite collision graph in auxiliary variables instead of
eliminating all intermediate states. Positions, unit normals, incoming
and outgoing unit velocities, and flight lengths through $m$ collisions
require at most $a(m+1)$ real variables. The flight equations, circle
equations and reflections have bounded degree. Positivity of flight
lengths and incidence, and exclusion of an earlier hit, are finite
Boolean combinations of bounded-degree sign conditions for each flight.
For example, the minimum squared distance of a straight segment from a
candidate center is tested at its endpoints and, when its orthogonal
projection lies inside the segment, at that projection. Positive
length denominators are cleared with their signs retained. The uniform
finite horizon gives a fixed finite candidate list per flight.
Consequently at most $a(m+1)$ polynomials of degree at most $d_*$
describe the graph, for fixed constants $a,d_*$. The same format
bounds hold after imposing a roof value, a coordinate line, or an
endpoint-gradient direction. Chart denominators are positive.

For precision, the bound used on sign-condition components is the
finite inequality in \cite[Section 3.2, equation (3.3)]{BPR}, with
$i=0$ and ambient variety $\R^r$:
\begin{equation}\label{eq:v44-sign-bound}
 b_0\le d(2d-1)^{r-1}\sum_{j=0}^r4^j\binom{s}{j}
       \le d(2d-1)^{r-1}5^s.
\end{equation}
Adding a bounding ball changes $s$ by one. A Boolean union has no
more components than the sum over its sign conditions. Here $r,s$
are linear in $m$, and $d$ is fixed, so the bound is exponential in
$m$. It is independent of polynomial coefficients, including $R$ and
the accumulated lattice translations. The fixed-dimension asymptotic
in the main theorem of that reference is not used with a varying dimension.
A projection of a connected set is connected. Thus the same component
bound controls every finite projection fiber of the physical graph.

For a regular one-dimensional level, almost every horizontal or
vertical line has a finite intersection with it. These intersections
are bounded by \eqref{eq:v44-sign-bound}, applied before projection.
Integrating the two coordinate multiplicities bounds respectively
$\int|dx|$ and $\int|dp|$. Since arclength is at most
$|dx|+|dp|$, and each coordinate range has length two, this proves
the length estimate. Straight segments parallel to a slicing line
cause an exceptional value only, and are counted by the other slicing.

For the endpoint assertion orient a regular level by its unit normal
$N=\nabla f/|\nabla f|$. If $T$ is its unit tangent, its curvature is
\begin{equation}\label{eq:v44-gauss-curvature}
 k_f=\frac{T^{\mathsf T}\nabla^2 f\,T}{|\nabla f|}>0.
\end{equation}
Thus its Gauss map has no constant interval on the near-normal portion.
A fiber of this map is semialgebraic: by the first variation its
unnormalized normal is
$R(-a(x)p_0,a(y)p_m)$. Parallelism to a specified unit direction and
the positive dot product are bounded-degree conditions on the same
auxiliary graph. Each fiber is finite, with the component bound just
proved. The one-dimensional area formula for the Gauss map gives
$\int |k_f|\,ds\le2\pi C_1C_2^m$. Adjust $C_1$.

One may first apply both area formulas on compact regular exhaustions.
Their right sides do not depend on the exhaustion. Monotone convergence
therefore includes portions approaching grazing, chart or decision
boundaries. Exceptional roof levels are discarded as permitted by the
coarea formula. No boundary-neighborhood mass is converted to a density
bound in this argument.
\end{proof}

\begin{theorem}[Complete fixed-count height bound]
\label{thm:v44-complete-density-bound}
There are $C>0$ and $A>1$, independent of $m,R$, such that
\begin{equation}\label{eq:v44-full-roof-height}
 \|\rho_{m,R}\|_\infty\le CA^m\qquad(m\ge1,
                                      R\in[0.45,0.47]).
\end{equation}
For any family of bounded measurable insertions $w_{n,R}$, all of
supremum norm at most $M$, the complete actual-return components obey
\begin{equation}\label{eq:v44-weighted-height}
 \operatorname*{ess\,sup}_{t}
 \sum_{n=1}^{m}\sum_{k\in\Z^2}
       |p^w_{n,R}(k,m,t)|\le MCA^m/c.
\end{equation}
In particular a prescribed component satisfies this bound with no
regularity condition on its insertion. No source guard or count
cutoff changes the measure in these statements.
\end{theorem}
\begin{proof}
Work first on one regular word in an initial rational chart. If
$|p_0|\ge\delta_0$, the nonstationary estimate
\eqref{eq:uniform-roof-gradient} gives a fixed positive lower gradient
bound. If $|p_0|<\delta_0$ and $|p_m|\ge\delta_0$, use instead
$|\partial_p L_m|=R|p_m B_m|\ge R_-\delta_0v_0^m$.
The rational change of the initial angle has derivative at least one.
The gradient is therefore at least $c_0v_0^m$ on the union of these
two regions, for a fixed $c_0>0$. The initial collision density is at
most $1/(2\pi)$ in this chart. Coarea and the level-length estimate
give a bound $C(C_2/v_0)^m$ for this part of the roof density.

On the remaining near-normal part change to the two endpoint charts.
The Jacobian weight is at most $v_0^{-m}/\pi$.
Equations~\eqref{eq:v44-chart-convexity} and
\eqref{eq:v44-gauss-curvature} give
\[
 \int_{\{f=t\}}\frac{ds}{|\nabla f|}
       \le\lambda_0^{-1}\int_{\{f=t\}}|k_f|\,ds
       \le CC_2^m,
\]
where both integrals are restricted to the near-normal part. This
argument does not extend the endpoint domain to its convex hull.
Coarea again gives the same exponential bound. Critical points on a
regular word are isolated by Proposition~\ref{prop:general-edge};
the critical roof levels do not affect this almost-everywhere bound.
The regular exhaustions in the preceding lemma include every portion
of each open physical branch. Singular source states themselves have
zero collision measure.

There are at most $D^m$ center words for a fixed finite alphabet of
size $D$, and only a fixed number of initial and endpoint charts.
Absorbing these factors in $A$ proves \eqref{eq:v44-full-roof-height}.
The constants are determined by the one-flight format, not by a
prepared singularity neighborhood or a lower interior incidence.

Finally, on $Y_R^*$ the events
$\{K_{n,R}=k,N_{n,R}=m\}$ are disjoint as $(n,k)$ varies with $m$
fixed: at collision $m$ a trajectory has a unique return index and
lattice displacement, when it returns at all. Their roofs are all
$L_{m,R}$. Pushforward domination and
$\nu_R^*\le c^{-1}\nu$ imply the measure inequality
$\sum_{n,k}|\mu^w_{n,R}(k,m,\cdot)|\le (M/c)(L_{m,R})_*\nu$.
Radon--Nikodym derivatives prove \eqref{eq:v44-weighted-height}.
\end{proof}

\paragraph{Interpretation.}
The exponential bound supplies uniform height control for the complete
finite-collision source and all its bounded restrictions. It does not
have the $m^{-2}$ scale of the arithmetic local law. Its important
consequence for raw extraction is that divergent density germs cannot
occur in these physical components, even though they occur in the
larger constructible $L^1$ class allowed in the earlier preparation lemma.
